% 图片文档放在source目录下面
% 2018年3月21日17:09:16 @ SWJTU by GeosciMan v1.0

\documentclass[12pt]{ctexbeamer}	%声明文档类型为beamer幻灯片
\mode <presentation>
\usetheme{Madrid}%%Warsaw [secheader]
%\usecolortheme{beaver}
\beamersetaveragebackground{black!10}

\setbeamercovered{transparent}
%\setbeamertemplate{navigation symbols}{}

\usefonttheme{professionalfonts}
\useinnertheme{circles}%{rectangles}
\setbeamertemplate{itemize item}{$\circledast$}%{\checkmark}


\setbeamertemplate{title page}
{
%  \vbox{}
  \vfill
  \begin{centering}
  \includegraphics[height=2.0cm]{source/swjtu_logo_fig_new.eps}
        \vskip0.6em\par%
    \begin{beamercolorbox}[sep=8pt,center,shadow=true,rounded=true]{title}
      \usebeamerfont{title}\inserttitle\par%
      \ifx\insertsubtitle\@empty%
      \else%
        \vskip0.25em%
        {\usebeamerfont{subtitle}\usebeamercolor[fg]{subtitle}\insertsubtitle\par}%
      \fi%
    \end{beamercolorbox}%
    \vskip1em\par
    \begin{beamercolorbox}[sep=8pt,center]{author}
      \usebeamerfont{author}
 %     \insertauthor
      {
        \begin{tabular}{cc}
  汇报人：&\insertauthor\\
  导\quad 师：& \advisor
  \end{tabular}
      }
    \end{beamercolorbox}
    \begin{beamercolorbox}[sep=8pt,center]{institute}
      \usebeamerfont{institute}\insertinstitute
    \end{beamercolorbox}
    \begin{beamercolorbox}[sep=8pt,center]{date}
      \usebeamerfont{date}\insertdate
    \end{beamercolorbox}\vskip0.5em
  %  {\usebeamercolor[fg]{titlegraphic}\inserttitlegraphic\par}
  \end{centering}
  \vfill
}

\hypersetup{pdfpagemode=FullScreen} % makes your presentation go automatically to full screen

\usepackage{fontspec,xunicode,xltxtra}

\usepackage{multimedia}	%让文档支持多媒体
\usepackage{graphics}	%让文档支持图片
\usepackage{hyperref}	%让文档支持超链接
\usepackage{booktabs}	%让文档支持三线表格
\usepackage{amsmath}	%ams可以让文档支持数学公式
\usepackage{amsfonts}
\usepackage{amssymb}
\usepackage{color}
\usepackage{graphicx,psfrag}
\usepackage{epsfig}

\title{季节性形变约束软流圈的瞬态粘弹性\\阅读报告}	%幻灯片标题
\author[陈昌福]{陈昌福}
\def\advisor{袁林果教授}

\institute[西南交通大学]{
大地测量学\\
西南交通大学，测绘遥感信息系
}	%作者单位
\date{\today}

\setcounter{tocdepth}{1}

\begin{document}
\begin{frame}
\titlepage
\end{frame}

\AtBeginSection[]{
  \frame{
    \frametitle{目录}
    \tableofcontents[currentsection]
  }
}

\begin{frame}
\frametitle{目录}
\tableofcontents
\end{frame}

%%%%%%%%%%%%%%%%%%%%%%%%%%%%%%%%%%%%%%%%%%%%%%%%%%%%%%%%%%%%%%%%%%%%%%%%%%%%%%%

\section{研究背景}
\begin{frame}
\frametitle{研究背景}

\begin{itemize}
    \item<1-> GPS监测地震周期的表征模式，区分地震形变的不同阶段
    \item<2-> 时空依赖的震后形变的本构关系(constitutive laws)研究仍具有较大的争议
    \item<3-> 时间跨度不同，预测值和观测值存在差异
    \item<4-> 负荷形变的研究和震后形变研究存在一定的差异
\end{itemize}
\end{frame}

%%%%%%%%%%%%%%%%%%%%%%%%%%%%%%%%%%%%%%%%%%%%%%%%%%%%%%%%%%%%%%%%%%%%%%%%%%%%%%%

\section{理论解释}
\begin{frame}
\frametitle{机制解释}
\begin{block}{震后余滑}
	\begin{itemize}
        \item 破裂断层或其附近迅速衰减滑动
		\item 适合于近场观测的解释
	\end{itemize}
\end{block}
\pause
\begin{block}{粘弹性延迟}
	\begin{itemize}
		\item 同震应力的粘弹性松弛
        \item 多个松弛时间以及应力依赖的粘性
	\end{itemize}
\end{block}
\pause
\begin{center}
    \alert {理论解释对震后形变的各自贡献比较难于区分
    \\但对解释地震周期以及余震和震后形变的能量联系较为关键}
\end{center}
\end{frame}

%%%%%%%%%%%%%%%%%%%%%%%%%%%%%%%%%%%%%%%%%%%%%%%%%%%%%%%%%%%%%%%%%%%%%%%%%%%%%%%

\section{研究方法}

\begin{frame}
\frametitle{数据}
\begin{itemize}
    \item<1-> GRACE数据
        \begin{itemize}
            \item 加回负荷部分的理论值
            \item 经验改正大地震引起的重力值变化
        \end{itemize}
    \item<2-> GPS数据
        \begin{itemize}
            \item 提取非构造信号
            \item 加权平均得到GRACE的采样率
        \end{itemize}
\end{itemize}
\end{frame}

\begin{frame}
\frametitle{结果讨论}
\begin{itemize}
    \item<1-> 对粘度提供独立的约束，解释早期的震后形变
    \item<2-> 在周年左右的时间尺度上，地球的弹性结构仍然独立潮汐研究(GNSS)的质疑 
    \item<3-> 周年负荷信号导出的值和长时间序列的结果一致
    \item<4-> 两种理论解释机制是相互作用？震后余滑可认为是释放部分同震应力？
\end{itemize}
\end{frame}

%%%%%%%%%%%%%%%%%%%%%%%%%%%%%%%%%%%%%%%%%%%%%%%%%%%%%%%%%%%%%%%%%%%%%%%%%%%%%%%

\section{总结评论}

\begin{frame}
\frametitle{总结评论}
\begin{itemize}
    \item<1-> 由于构造应力，没有考虑应变特性的非线性性质
    \item<2-> 相同时空尺度下，负荷形变的研究结果和震后形变的研究结果不兼容
    \item<3-> 震后松弛形变贡献的上界，未给出两者具体差异的原因
    \item<4-> 地震频段及潮汐频段的耗散研究结果相兼容
    \item<5-> 两种理论机制可综合应用，单个应用可能会存在某些不兼容
\end{itemize}
\end{frame}

%%%%%%%%%%%%%%%%%%%%%%%%%%%%%%%%%%%%%%%%%%%%%%%%%%%%%%%%%%%%%%%%%%%%%%%%%%%%%%%


\begin{frame}
\centerline{\Large 欢迎大家批评指正，谢谢！}
\end{frame}

\end{document}
